\chapter{2015 Nobel Prize in Economics}
\textbf{Laureates: }Angus Deaton (UK/USA)

\section*{Reason for Award}
For his analysis of consumption, poverty, and welfare

\section{What They Did}
Angus Deaton made foundational contributions to the analysis of consumption,
poverty, and welfare. Born in Edinburgh, Scotland, in 1945, Deaton earned his
doctorate from the University of Cambridge before becoming a professor at
Princeton University.

He developed the almost ideal demand system (AIDS), a flexible functional form
for estimating consumer demand that became standard in empirical industrial
organization and policy analysis. The AIDS model allows researchers to estimate
how consumers respond to price changes and income variations while satisfying
theoretical restrictions like symmetry and adding-up. This framework has been
used to evaluate tax policies, analyze market structures, and assess the
distributional effects of economic reforms.

Deaton's work on household survey methodology improved how economists measure
poverty and living standards in developing countries. He showed why consumption
or expenditure measures can be useful for assessing material welfare, particularly
where income is volatile or difficult to measure accurately, while also emphasizing
the measurement problems that affect both approaches. His research on survey
design and measurement error has improved the quality of data used in poverty
analysis.

He analyzed the relationship between consumption, income, and wealth, showing
how income expectations and constraints affect current consumption. Deaton's work
on consumption smoothing examined how households use saving, borrowing, and other
resources to manage risk, while recognizing that credit constraints and incomplete
insurance can prevent consumption from remaining stable.

His broader research also examined links among health, consumption, and economic
well-being. These connections show why measures of living standards cannot be
separated entirely from health and mortality, and why economic resources and
health outcomes may influence one another.

Deaton also studied development policy and the measurement of poverty and welfare.
Research conducted with collaborators, including work on South Africa, examined
relationships among inequality, health, and economic conditions, illustrating how
household data can illuminate the effects of political and economic institutions.

\section{Impact on Economic Thinking}
Deaton transformed consumption theory and poverty measurement. The almost ideal
demand system provided researchers with practical tools for analyzing consumer
behavior and evaluating policy impacts. His work on household surveys established
standards for collecting and analyzing microdata in developing countries.

Deaton's analysis of consumption smoothing showed how households manage risk
through saving, borrowing, and informal support, while also revealing the limits
created by credit constraints and incomplete insurance. His broader work connecting
individual well-being to aggregate outcomes influenced research on health, poverty,
and public policy.

Deaton's later work examined inequality, development, and the measurement of
human progress. His book ``The Great Escape'' (2013) traced how improvements in
health and living standards have transformed human welfare while also emphasizing
the unequal distribution of those gains and the dangers of judging progress by
GDP alone.

His interdisciplinary approach enriched economics, demography, and public health.
He showed how economic methods can address fundamental questions about human
well-being, bridging the gap between quantitative analysis and qualitative
understanding of human experience.

Deaton's work on measurement has influenced international poverty research and
the design of social-protection programs. His emphasis on high-quality microdata
has raised standards for empirical research in development economics.

\section{Modern Perspectives Today}
Deaton's methodological contributions remain important in empirical economics.
The almost ideal demand system is widely used in consumer research, policy
evaluation, and industrial organization. His work on poverty measurement informs
international development research and policy analysis.

Research on consumption smoothing continues to explore how financial development,
social insurance, and credit constraints affect household welfare. Deaton's work
on health and living-standard measurement informs research on health inequality
and the relationship between health and economic outcomes. These measurement
insights are also relevant to evaluating public-health interventions, including
responses to major health crises.

His work on inequality and development has contributed to debates about the
relationship between economic resources, political institutions, and human
welfare. Policies addressing inequality require understanding both economic
mechanisms and the institutions through which resources and opportunities are
distributed.

As global inequality and health disparities persist, Deaton's analytical tools
remain essential for understanding and addressing these challenges. His emphasis
on measurement quality and interdisciplinary approaches continues to guide
research on poverty, health, and well-being in the twenty-first century.

\subsection*{Key References}
\begin{itemize}
    \item Deaton, A., \& Muellbauer, J. (1980). ``An Almost Ideal Demand System.'' \textit{American Economic Review}, 70(3), 312--326.
    \item Deaton, A. (1997). \textit{The Analysis of Household Surveys: A Microeconometric Approach to Development Policy}. Johns Hopkins University Press.
\end{itemize}
